\begin{annexitem}{annex:p17.rat8.3.binary-reference}{thm:p17-pt-binary-reference}{Binary risk attitudes and fixed-final-outcome reversal}
\begin{proof}
For \(x>0\), \(v'(x)=\alpha x^{\alpha-1}>0\) and
\(v''(x)=\alpha(\alpha-1)x^{\alpha-2}<0\). For \(x<0\),
\(v'(x)=\lambda\beta(-x)^{\beta-1}>0\) and
\(v''(x)=-\lambda\beta(\beta-1)(-x)^{\beta-2}>0\).
Continuity and \(v(0)=0\) extend monotonicity and the one-sided
concavity/convexity to their closed half-lines; no differentiability
at the kink is required.

The binary gain lottery has value \(w^+(p)b^\alpha\); its sure
mean has value \((pb)^\alpha=p^\alpha b^\alpha\). Division by
positive \(b^\alpha\) gives the stated gain condition and its tie
case. The binary loss lottery has value
\(-\lambda w^-(p)b^\beta\); its sure mean has value
\(-\lambda p^\beta b^\beta\). Division by negative
\(-\lambda b^\beta\) reverses the inequality and gives the
loss condition. No assertion about the pattern follows without
matching each weighting function against these thresholds.

At reference zero and with the stipulated identity weighting,
the certain-4 score is \(\sqrt4=2\), while the lottery score is
\(\tfrac12\sqrt{10}\). Since \(10<16\), the latter is less
than 2. At reference 10, the certain outcome is deviation \(-6\)
with score \(-2\sqrt6\). The risky deviations are 0 and \(-10\),
with score \(-\sqrt{10}\). Since \(10<24\),
\(-\sqrt{10}>-2\sqrt6\), so the lottery is preferred.
These are the same final-consequence distributions in both
comparisons. For a fixed final-consequence utility \(u\), their
expected-utility scores remain \(u(4)\) and
\(\tfrac12u(10)+\tfrac12u(0)\), independently of the label
\(r\). Thus an expected-utility reversal would require changing
some represented input or utility, not merely relabeling a reference.
\end{proof}
\end{annexitem}

\begin{annexitem}{annex:p17.rat8.3.source-calculations}{thm:p17-pt-source-calculations}{Finite monetary scores and the source's St.~Petersburg convention}
\begin{proof}
The expected-money scores are \(4\) and
\(\tfrac12\cdot10+\tfrac12\cdot0=5\).
For power utility with \(a>0\), the utility scores are
\(4^a\) and \(\tfrac12 10^a\). The first exceeds the
second precisely when \(2>(5/2)^a\). Since the logarithm
is strictly increasing and \(\log(5/2)>0\), this is
\(a<\log2/\log(5/2)\). Equality gives indifference.
At \(a=.7\), direct exponentiation gives approximately
\(4^{.7}=2.639015822\) and
\(\tfrac12 10^{.7}=2.505936168\), verifying the source's
rounded scores.

With independent fair flips,
\(\Pr(K=k)=2^{-(k+1)}\) for every \(k\geq0\).
The probability of tails forever is zero because its
probability is bounded by \(2^{-n}\) for every \(n\).
The assigned zero payoff on that null event makes the payoff
defined everywhere and adds nothing to either expectation.
The expected contribution of each \(k\geq1\) is
\(2^{k-1}2^{-(k+1)}=1/4\), while \(k=0\) contributes
zero. Nonnegative partial sums therefore grow as \(n/4\),
so the monetary expectation is infinite. With power utility,
\[
 \sum_{k=1}^\infty
       2^{a(k-1)}2^{-(k+1)}
   =\frac14\sum_{j=0}^\infty 2^{(a-1)j}
   =\frac{1}{4(1-2^{a-1})}.
\]
The geometric ratio lies strictly between zero and one for
\(0<a<1\), proving convergence. This calculation uses the
source's zero-on-immediate-head payoff, not the differently
indexed version paying a positive amount on the first head.
\end{proof}
\end{annexitem}

\begin{annexitem}{annex:p17.rat8.3.coalescing}{thm:p17-pt-coalescing}{Telescoping masses and presentation invariance}
\begin{proof}
Cumulative probabilities are nested, so every difference of
their values under an increasing \(w^+\) or \(w^-\) is
nonnegative. Consecutive differences telescope. For losses,
the increasing cumulative sequence begins at zero and ends at
\(P_-\), giving total \(w^-(P_-)-w^-(0)=w^-(P_-)\).
For gains, the decreasing upper-tail sequence begins at
\(P_+\) and ends at zero, giving total
\(w^+(P_+)-w^+(0)=w^+(P_+)\). If all mass is strictly
on one side, the sum is \(w^\pm(1)=1\).

For a half-gain/half-loss prospect with no zero mass, take
strictly increasing piecewise-linear weighting functions
through \((0,0),(.5,.4),(1,1)\), the same on both signs.
Their combined mass is .8. Replacing the middle value by
.6 makes it 1.2, still with valid increasing normalized
functions. Hence there is no general global mass-one
conclusion for mixed prospects. Zero outcomes contribute
zero value; an arbitrary convention for their unused
weight does not alter these nonzero contributions.

Combining equal-outcome branches produces exactly the
same probability attached to each distinct deviation
before and after a split. Every lower and upper cumulative
probability in the definition is therefore identical.
The value and score sums are identical too. Equivalently,
if tied branches are retained adjacently, their weights
telescope within the tie and multiply the same value.
Thus coalescing does not rely on a particular fit or
curvature. This proves the impossibility of a pure
presentation reversal for a fixed CPT representation of
the unchanged distributions.

For the displayed three-gain lottery, the probabilities
in decreasing outcome order are .2, .4, .4, so cumulative
upper tails are .2, .6, 1. The weights are exactly the
three stated differences. The stipulated piecewise-linear
function is well defined and strictly increasing because
both coordinates increase strictly at each consecutive
node. Its differences are .30, .55-.30=.25, and
1-.55=.45. Their sum is one, but equality of the last
two branch probabilities does not force equality of
their cumulative weights.
\end{proof}
\end{annexitem}

\begin{annexitem}{annex:p17.rat8.3.sampling}{thm:p17-pt-sampling}{Exact sampling masses and a small-sample countercase}
\begin{proof}
With independent indicators of probability \(1/5\), each
particular sequence containing \(j\) rare outcomes has
probability \((1/5)^j(4/5)^{10-j}\). There are
\(\binom{10}{j}\) such sequences, so their disjoint
sum gives the binomial mass function. Summing the
indicator expectations gives \(\mathbb E[N]=10/5=2\).
Consequently
\[
 \Pr(N<2)=(4/5)^{10}+10(1/5)(4/5)^9
                =.3758096384,
\]
and
\[
 \Pr(N=2)=45(1/5)^2(4/5)^8=.3019898880.
\]
Complementation gives
\(\Pr(N>2)=1-.3758096384-.3019898880=.3222004736\).
The three disjoint masses sum exactly to one. The
chapter's two rounded values compare below-mean and
above-mean counts while omitting the substantial
equal-mean mass; they are not complementary events.

For four independent samples the mean is \(4/5=.8\).
A count is under that mean exactly when it is zero,
and over it exactly when it is at least one. These
events have probabilities \((4/5)^4=.4096\) and
\(1-(4/5)^4=.5904\), respectively. Thus even at
the same .2 rare-outcome probability, changing sample
size reverses the direction of the comparison. Skewness
alone is insufficient for the unrestricted verbal claim.
This is an analytical counterexample, not a measured
description--experience response.
\end{proof}
\end{annexitem}
